% preamble-exam-zh.tex — 共享 preamble
% 用于所有 exam-zh 模板
%
% 样式策略：
%   屏幕 → 语义彩色（蓝=关键想法，红=易错，绿=路线，灰=提示/教师备注）
%   打印 → 自动降级为黑白（通过 print mode 或 monochrome 选项）

\documentclass{exam-zh}

% ---------- 基础配置 ----------
\examsetup{
  page/size = a4paper,
  page/show-head = false,
  page/show-foot = false,
  sealline/show = false,
  font = times,
  math-font = xits,
}

\usepackage{tcolorbox}
\usepackage{needspace}
\tcbuselibrary{skins,breakable}

% ---------- 打印/屏幕颜色切换 ----------
% 用法：\documentclass[monochrome]{exam-zh} 即可切换为纯黑白
% 注意：不要使用 \if@xxx 放在 makeatother 外部；这里改成普通条件名。
\newif\ifeduprintcolor
\eduprintcolortrue

\makeatletter
\@ifclasswith{exam-zh}{monochrome}{\eduprintcolorfalse}{}
\makeatother

% ---------- 语义颜色定义 ----------
% 屏幕：有色彩区分；打印：降级为灰度
\ifeduprintcolor
  % --- 屏幕彩色模式 ---
  \definecolor{edu-blue}{RGB}{47, 82, 122}       % 关键想法
  \definecolor{edu-blue-bg}{RGB}{243, 246, 252}
  \definecolor{edu-red}{RGB}{180, 40, 40}         % 易错提醒
  \definecolor{edu-red-bg}{RGB}{255, 245, 240}
  \definecolor{edu-green}{RGB}{34, 100, 60}       % 路线图
  \definecolor{edu-green-bg}{RGB}{242, 250, 245}
  \definecolor{edu-orange}{RGB}{160, 90, 0}       % 训练目标
  \definecolor{edu-orange-bg}{RGB}{255, 248, 235}
  \definecolor{edu-gray}{RGB}{100, 100, 100}      % 提示/教师备注
  \definecolor{edu-gray-bg}{RGB}{248, 248, 248}
  \definecolor{edu-purple}{RGB}{90, 50, 120}      % 思考题
  \definecolor{edu-purple-bg}{RGB}{248, 244, 252}
\else
  % --- 打印黑白模式 ---
  \definecolor{edu-blue}{gray}{0.15}
  \definecolor{edu-blue-bg}{gray}{0.96}
  \definecolor{edu-red}{gray}{0.25}
  \definecolor{edu-red-bg}{gray}{0.97}
  \definecolor{edu-green}{gray}{0.20}
  \definecolor{edu-green-bg}{gray}{0.97}
  \definecolor{edu-orange}{gray}{0.25}
  \definecolor{edu-orange-bg}{gray}{0.98}
  \definecolor{edu-gray}{gray}{0.40}
  \definecolor{edu-gray-bg}{gray}{0.97}
  \definecolor{edu-purple}{gray}{0.30}
  \definecolor{edu-purple-bg}{gray}{0.98}
\fi

% ---------- 自定义教学环境 ----------
% 共性：enhanced + breakable（跨页不断裂）

% 原题卡片 — 强边框，突出显示
\newtcolorbox{problemcard}{
  enhanced, breakable,
  colback=edu-blue-bg,
  colframe=edu-blue,
  boxrule=1.2pt,
  arc=0pt,
  left=3mm, right=3mm, top=3mm, bottom=3mm,
}

% 解题路线图 — 左边线 + 浅底
\newtcolorbox{routebox}{
  enhanced, breakable,
  colback=edu-green-bg,
  colframe=edu-green!30,
  boxrule=0pt,
  borderline west={2.5pt}{0pt}{edu-green},
  left=3mm, right=3mm, top=3mm, bottom=3mm,
}

% 关键想法 — 蓝色左边线
\newtcolorbox{keyideabox}{
  enhanced, breakable,
  colback=edu-blue-bg,
  colframe=edu-blue!30,
  boxrule=0pt,
  borderline west={2.5pt}{0pt}{edu-blue},
  left=3mm, right=3mm, top=3mm, bottom=3mm,
}

% 易错提醒 — 红色左边线
\newtcolorbox{mistakebox}{
  enhanced, breakable,
  colback=edu-red-bg,
  colframe=edu-red!20,
  boxrule=0pt,
  borderline west={3pt}{0pt}{edu-red},
  left=3mm, right=3mm, top=2.5mm, bottom=2.5mm,
}

% 提示 — 灰色虚线左边线
\newtcolorbox{hintbox}{
  enhanced, breakable,
  colback=edu-gray-bg,
  colframe=edu-gray!20,
  boxrule=0pt,
  borderline west={2pt}{0pt}{edu-gray, dashed},
  left=3mm, right=3mm, top=2mm, bottom=2mm,
}

% 思考题 — 紫色虚线左边线
\newtcolorbox{thinkbox}{
  enhanced, breakable,
  colback=edu-purple-bg,
  colframe=edu-purple!20,
  boxrule=0pt,
  borderline west={2pt}{0pt}{edu-purple, dashed},
  left=2.5mm, right=2.5mm, top=2.5mm, bottom=2.5mm,
}

% 教师备注 — 灰色左边线，小字
\newtcolorbox{teachernote}{
  enhanced, breakable,
  colback=edu-gray-bg,
  colframe=edu-gray!15,
  boxrule=0pt,
  borderline west={2pt}{0pt}{edu-gray!60},
  left=2.5mm, right=2.5mm, top=2.5mm, bottom=2.5mm,
  fontupper=\small,
  colupper=black!60,
}

% 训练目标 — 橙色左边线，小字
\newtcolorbox{traininggoal}{
  enhanced, breakable,
  colback=edu-orange-bg,
  colframe=edu-orange!15,
  boxrule=0pt,
  borderline west={1.5pt}{0pt}{edu-orange},
  left=2mm, right=2mm, top=1mm, bottom=1mm,
  fontupper=\small,
}

% 答题区 — 无边框，纯留白
\newcommand{\answerarea}[1][40mm]{%
  \vspace{2mm}%
  \begin{tcolorbox}[
    enhanced, breakable,
    colback=white,
    colframe=white,
    boxrule=0pt,
    height=#1,
    valign=top,
    bottom=0mm,
    after skip=0mm,
  ]
  \end{tcolorbox}%
}

% ---------- 字体设置 ----------
% exam-zh (ctexbook) already sets CJK fonts via ctex.
% Only override if specific fonts are needed.
% macOS: Songti SC, STHeiti, STFangsong
% Linux: SimSun, SimHei, FangSong

% ---------- 页面设置 ----------
% exam-zh (ctexbook) already loads geometry.
% Override margins if needed:
% \geometry{top=16mm, bottom=16mm, left=14mm, right=14mm}

\examsetup{
  question/show-answer = true,
  fillin/show-answer = true,
  solution/show-solution = show-stay,
}

% ---------- 教师版解答样式 ----------
\newtcolorbox{solutionblock}{
  enhanced, breakable,
  colback=edu-blue-bg,
  colframe=edu-blue!25,
  boxrule=0pt,
  borderline west={2pt}{0pt}{edu-blue!50},
  arc=0pt,
  left=3mm, right=3mm, top=2mm, bottom=2mm,
  before skip=2mm,
  after skip=3mm,
  fontupper=\small,
}

% 解答步骤编号
\newcounter{solstep}
\newcommand{\solstep}[2]{%
  \stepcounter{solstep}%
  \par\medskip
  \noindent\hangindent=6mm
  {\color{edu-blue}\bfseries\textcircled{\scriptsize\thesolstep}\enspace #1}\enspace #2\par
}

\begin{document}

\section*{三垂直与旋转等腰直角专题（一） · 练习卷（教师版）}

\section*{一、选择题（每题 12 分，共 48 分）}


\begin{keyideabox}
  \textbf{核心模型：一线三垂直}
  已知等腰直角 $\triangle ABC$ 中 $\angle BAC = 90^\circ$，$AB = AC$。过直角顶点 $A$ 作一条直线，从另外两个顶点 $B, C$ 向该直线作垂线，垂足为 $B_1, C_1$。由直角互余可证 $\triangle ABB_1 \cong \triangle CC_1A$ (AAS)。水平位移与竖直位移发生互换。
\end{keyideabox}



\begin{question}[points=12]
  直线 $y = -x + 4$ 与 $x$ 轴、 $y$ 轴分别交于点 $A$ 和点 $B$。以 $AB$ 为直角边在第一象限作等腰 Rt$\triangle ABC$ （直角顶点为 $A$，即 $\angle BAC = 90^\circ$），则点 $C$ 的坐标为
  \begin{choices}
    \item $(8, 4)$
    \item $(4, 8)$
    \item $(4, 4)$
    \item $(8, 0)$
  \end{choices}
  \paren[A]
  \begin{solutionblock}
    $A(4,0)$，$B(0,4)$。做辅助线 $x$ 轴，$BB_1 \perp x$ 轴于 $B_1(0,0)$，$CC_1 \perp x$ 轴于 $C_1$。$\triangle ABB_1 \cong \triangle CC_1A$ $\implies CC_1 = AB_1 = 4 \implies y_c = 4$；$AC_1 = BB_1 = 4 \implies x_c = x_a + AC_1 = 4+4=8$。坐标为 $(8, 4)$。
  \end{solutionblock}
\end{question}



\begin{question}[points=12]
  一线三垂直模型的核心是证全等，其全等的判定依据是
  \begin{choices}
    \item AAS
    \item SAS
    \item SSS
    \item HL
  \end{choices}
  \paren[A]
  \begin{solutionblock}
    通过三直角证明两锐角相等（余角性质），以及腰相等，构成角角边 (AAS) 证明全等。
  \end{solutionblock}
\end{question}



\begin{question}[points=12]
  直线 $y = -2x + 6$ 与 $x$ 轴、 $y$ 轴分别交于点 $A$ 和点 $B$。以 $AB$ 为直角边在第一象限作等腰 Rt$\triangle ABC$ （直角顶点为 $A$，$\angle BAC = 90^\circ$），则点 $C$ 的坐标为
  \begin{choices}
    \item $(9, 3)$
    \item $(6, 3)$
    \item $(3, 6)$
    \item $(9, 6)$
  \end{choices}
  \paren[A]
  \begin{solutionblock}
    $A(3,0)$，$B(0,6)$。作 $x$ 轴垂线，$B_1(0,0)$，垂线段 $BB_1 = 6$，$AB_1 = 3$。\\ $\triangle ABB_1 \cong \triangle CC_1A \implies CC_1 = AB_1 = 3 \implies y_c = 3$；$AC_1 = BB_1 = 6 \implies x_c = x_a + 6 = 9$。坐标为 $(9, 3)$。
  \end{solutionblock}
\end{question}



\begin{question}[points=12]
  已知点 $A(2, 0)$，点 $B(0, 2)$。若以 $AB$ 为直角边在第一象限作等腰 Rt$\triangle ABC$ （直角顶点为 $A$，$\angle BAC = 90^\circ$），则点 $C$ 的坐标为
  \begin{choices}
    \item $(4, 2)$
    \item $(2, 4)$
    \item $(4, 4)$
    \item $(2, 2)$
  \end{choices}
  \paren[A]
  \begin{solutionblock}
    $A(2,0)$，$B(0,2)$。同样地，$\triangle ABB_1 \cong \triangle CC_1A \implies CC_1 = AB_1 = 2 \implies y_c = 2$；$AC_1 = BB_1 = 2 \implies x_c = 2+2=4$。坐标为 $(4, 2)$。
  \end{solutionblock}
\end{question}


\section*{二、填空题（每题 12 分，共 12 分）}


\begin{question}[points=12]
  已知两点 $A(3, 0)$ 和 $B(0, 4)$。以 $AB$ 为直角边在第一象限作等腰 Rt$\triangle ABC$ （直角顶点为 $A$），则点 $C$ 的坐标为 \fillin。
  \paren[$(7, 3)$]
  \begin{solutionblock}
    $A(3,0)$，$B(0,4)$。$\triangle ABB_1 \cong \triangle CC_1A \implies CC_1 = AB_1 = 3 \implies y_c = 3$；$AC_1 = BB_1 = 4 \implies x_c = x_a + 4 = 7$。坐标为 $(7, 3)$。
  \end{solutionblock}
\end{question}


\section*{三、解答题（共 40 分）}

\needspace{8\baselineskip}

\begin{problem}[points=20]
  直线 $y = -\dfrac{\sqrt{3}}{3}x + 1$ 与 $x$ 轴、$y$ 轴分别交于点 $A$、点 $B$。以 $AB$ 为直角边在第一象限作等腰 Rt$\triangle ABC$ （直角顶点为 $A$，$\angle BAC = 90^\circ$）。

(1) 求点 $A$ 和点 $B$ 的坐标；
(2) 证明 $\triangle ABB_1 \cong \triangle CC_1A$ （其中 $B_1, C_1$ 为 $B, C$ 到 $x$ 轴的垂足），并求点 $C$ 的坐标。
  \answerarea[60mm]
  \setcounter{solstep}{0}
  \begin{solutionblock}
    \textbf{答案：}(1) $A(\sqrt{3}, 0)$，$B(0, 1)$；(2) $C(\sqrt{3}+1, \sqrt{3})$\par\medskip
    \solstep{求 A、B 坐标}{令 $y = 0 \implies x = \sqrt{3}$，故 $A(\sqrt{3}, 0)$。\\ 令 $x = 0 \implies y = 1$，故 $B(0, 1)$。}
    \solstep{证明 AAS 全等}{过直角顶点 $A$ 沿 $x$ 轴做辅助线。自 $B, C$ 向 $x$ 轴作垂线，垂足 $B_1(0,0)$，$C_1(x_c, 0)$。\\ $\angle AB_1B = \angle AC_1C = 90^\circ$。\\ $\because \angle B_1AB + \angle CAC_1 = 90^\circ$，且 $\angle B_1AB + \angle ABB_1 = 90^\circ$。\\ $\therefore \angle ABB_1 = \angle CAC_1$。\\ 在 $\triangle ABB_1$ 和 $\triangle CC_1A$ 中，有 $\angle AB_1B = \angle AC_1C = 90^\circ$，$\angle ABB_1 = \angle CAC_1$，且 $AB = AC$。\\ $\therefore \triangle ABB_1 \cong \triangle CC_1A$ (AAS)。}
    \solstep{利用对应边求 C 坐标}{由全等得 $CC_1 = AB_1 = \sqrt{3} \implies y_c = \sqrt{3}$。\\ $AC_1 = BB_1 = 1 \implies x_c = x_a + AC_1 = \sqrt{3} + 1$。\\ 故点 $C$ 的坐标为 $(\sqrt{3}+1, \sqrt{3})$。}
  \end{solutionblock}
\end{problem}


\needspace{8\baselineskip}

\begin{problem}[points=20]
  直线 $y = -x + 2$ 与 $x$ 轴、 $y$ 轴分别交于点 $A$ 和点 $B$。以 $AB$ 为直角边在第一象限内作等腰 Rt$\triangle ABC$ （直角顶点为 $A$，$\angle BAC = 90^\circ$）。

(1) 证明 $\triangle ABB_1 \cong \triangle CC_1A$ （$B_1, C_1$ 分别为 $B, C$ 到 $x$ 轴的垂足）；
(2) 求点 $C$ 的坐标并验证 $AC = AB$ 的长度。
  \answerarea[55mm]
  \setcounter{solstep}{0}
  \begin{solutionblock}
    \textbf{答案：}(1) AAS全等证明；(2) $C(4, 2)$，经过验证 $AB = AC = 2\sqrt{2}$\par\medskip
    \solstep{求两轴交点及全等证明}{$A(2, 0)$，$B(0, 2)$。\\ 证明同前：直角三角形互余 $\implies \angle ABB_1 = \angle CAC_1$。利用 AAS 判定两三角形全等。}
    \solstep{求 C 点坐标}{由全等得 $CC_1 = AB_1 = 2 \implies y_c = 2$。\\ $AC_1 = BB_1 = 2 \implies x_c = x_a + 2 = 2 + 2 = 4$。\\ 故点 $C$ 坐标为 $(4, 2)$。}
    \solstep{验证长度}{$AB = \sqrt{2^2 + 2^2} = \sqrt{8} = 2\sqrt{2}$。\\ $AC = \sqrt{(4-2)^2 + (2-0)^2} = \sqrt{2^2 + 2^2} = \sqrt{8} = 2\sqrt{2}$。\\ 长度相等，验证通过。}
  \end{solutionblock}
\end{problem}


\clearpage
\section*{参考答案}


\begin{keyideabox}
  \textbf{选择题}
  1. A  2. A  3. A  4. A
\end{keyideabox}



\begin{keyideabox}
  \textbf{填空题}
  1. (7, 3)
\end{keyideabox}



\begin{keyideabox}
  \textbf{解答题}
  第 1 题：(1) $A(\sqrt{3}, 0)$，$B(0, 1)$；(2) $C(\sqrt{3}+1, \sqrt{3})$
第 2 题：(1) AAS全等证明；(2) $C(4, 2)$，经过验证 $AB = AC = 2\sqrt{2}$
\end{keyideabox}



\end{document}
